\section{Equity, Diversity, and Inclusion}\label{equity-diversity-and-inclusion}

\stanceinfo{\textbf{CFES Congress 2019} [2019-01-04]}{2024-01-02}

\officialstance{The Canadian Federation of Engineering Students (CFES) believes that equity, diversity and inclusion (EDI) are essential to the sustainability of the engineering profession and its contribution to the public interest. Increased inclusion and participation of underrepresented groups bring different perspectives and experiences that help engineers when addressing problems within the engineering profession and throughout undergraduate curricula.}

\stanceheading{The Students' Position}
\begin{itemize}
 \item Increasing the diversity and minority group representation within engineering helps to tackle challenges and create solutions to engineering problems that benefit everyone.
 \item Engineering students should have access to equal opportunities in engineering regardless of their age, ethnicity, race, gender, sexual orientation, economic status, disability, religion, or other diverse backgrounds. These opportunities include access to resources and opportunities that further their growth as an engineer (i.e internships, engineering conferences, etc.).
 \item Increasing the diversity of students in engineering programs creates an enhanced engineering education by exposing students to a more significant amount of perspectives, experiences, and ways of thinking.
 \item Diversity will be fostered at a higher and more successful rate when in an inclusive environment by contributing to the retention of underrepresented groups. In order to create an inclusive environment, education and initiatives must be created around topics of diversity to educate individuals about increasing accessibility with diverse peoples and reduce unconscious bias.
 \item Indigenous inclusion in engineering is an important, yet under-discussed topic given that Canada is on traditional and unceded Indigenous lands and exists as a result of colonization; the effects of which are still present today.
 \item The CFES further endorses the Engineers Canada National Position Statement on Diversity and Inclusion, and the Engineers Canada 30 by 30 initiative on raising the percentage of newly licensed engineers who identify as women to 30 percent by the year 2030.
 \item Underrepresented groups are often at a disadvantage due to systemic barriers that disproportionately impact them, which can make it difficult to enter engineering and thrive.
 \item Creating an inclusive environment within engineering means to create a safe environment, culture and practice that welcomes and respects the contributions of individuals from diverse backgrounds, experiences and perspectives.
\end{itemize}

\stanceheading{Recommendations to Partners, Stakeholders, and Other Entities}
\begin{itemize}
 \item The CFES calls on the Engineers Canada ``30 by 30 Network'' to investigate the reasons why retention rates of women are low in undergraduate education and formally present to other stakeholders including the Engineers Canada Board of Directors.
 \item The CFES calls on Engineers Canada ``30 by 30 Network'' to begin developing secondary plans if 30 by 30 is not accomplished Canada wide, and integrate student voices in the new proposal.
 \item The CFES calls on Engineering Deans Canada (EDC) to continue promoting the recruitment of diverse identities within incoming engineering students with initiatives that promote equity within the engineering community. (i.e underrepresented community scholarships).
 \item The CFES calls on the Canadian Engineering Accreditation Board (CEAB) to integrate equity, diversity and inclusion (EDI) into the Graduate Attributes, similar to the impact of engineering on society and the environment, to describe the importance of EDI training and always ensuring EDI considerations are done in the engineering profession.
 \item The CFES calls on EDC to integrate an Associate Dean of Equity, Diversity and Inclusion or a position of similar standing within their institutions to achieve equitable, diverse and inclusive accessibility to engineering and provide support to underrepresented groups.
 \item The CFES calls on EDC to provide equity, diversity and inclusion training to their instructors and students, specifically students that represent the institution at large.
 \item The CFES calls on engineering faculties across Canada to hire diverse instructors that reflect the diversity of students they teach and ensure support for marginalized faculty.
 \item The CFES calls upon Institution-specific Administrators, with the support of Engineering Deans Canada to implement and enforce protocols to handle cases of discrimination and harassment within the engineering community.
 \item The CFES calls upon the CEAB and EDC to promote stewardship principles in curriculum, to ensure that students are better equipped to consider and champion EDI values in the workplace upon completion of their studies.
 \item Lastly, the CFES calls on our members to examine the engineering traditions that they uphold and consider with a critical lens whether or not they uphold and contribute to creating an inclusive environment for all. Furthermore, the CFES calls on members to seek out opportunities to promote EDI literacy amongst students, and create initiatives that support equity, diversity and inclusion in their institutions.
\end{itemize}

\stanceheading{What the CFES plans to do}

The CFES plans to strengthen our support and promotion of inclusion within the engineering student community by providing resources and sharing best practices to member schools regarding hosting inclusive events and developing inclusivity orientations/training within their respective communities.

\begin{itemize}
 \item Continue its inclusivity efforts by offering inclusivity orientations at all events, as well as further investigating event inclusivity risk assessment.
 \item Explore new partnerships with diversity and inclusion-focused organizations (Catalyst Canada, Actua, Women in Engineering, Canadian Mental Health Association, National Society of Black Engineers, Diversity Initiative at the Engineering Change Lab, etc).
 \item Prioritise the active incorporation of bilingualism within all CFES Activities, as well as within the core of the organization.
 \item Develop inclusive event planning guidelines specific to CFES activities to help identify, remove, and prevent barriers that may hinder anyone from fully participating in CFES Activities and ensure these considerations are taken from the early planning stages of each activity.
 \item Fulfill all the commitments set out in the Champions Charter to officially endorse Engineer's Canada's 30 by 30 initiative, including nominating a Champion at Spring Meeting each year that participates in the Champions' group and actively engages in the 30 by 30 network.
 \item Put into practise the Exclusionary \& Harmful Reporting Form that provides a space for feedback from members.
 \item Prioritize barrier-free events and activities at CFES conferences, with alternatives made available when this is not possible.
 \item Create a CFES Equity, Diversity and Inclusion Toolkit for members, that includes but is not limited to the following:
 \begin{itemize}
  \item Equity, diversity and inclusion training
  \item EDI certification program
  \item Best practices for external organizations
  \item Canadian statistics of engineering demographic information
  \item Access to CFES partners for local advocacy
  \item Feedback mechanisms
 \end{itemize}
\end{itemize}

\newpage
